\chapter{First contact}
\label{ch:contact}

\chapterbrief
  {Find the Pi and log into it.}
  {Hostnames, mDNS, DHCP --- how machines find each other.}
  {A shell prompt that says \code{martin@brewer}.}

\section{Log in}

\begin{shell}
# on your laptop
ssh martin@brewer.local
\end{shell}

The first time, you get this:

\begin{trace}
The authenticity of host 'brewer.local (192.168.1.42)' can't be established.
ED25519 key fingerprint is SHA256:xK3f...9dQ.
Are you sure you want to continue connecting (yes/no/[fingerprint])?
\end{trace}

Type \textbf{\code{yes}} and press Enter. Then your password (it shows nothing as
you type --- that is deliberate, not a broken keyboard).

\begin{trace}
martin@brewer:~ $
\end{trace}

That prompt means you are on the Pi. Everything you type now runs there.

\begin{concept}{What that fingerprint question is}
Every SSH server has its own identity key. The first time you connect, your
computer has no way to know whether this is really your Pi or something
pretending to be it, so it shows you the fingerprint and records your answer in
\code{\textasciitilde/.ssh/known\_hosts}.

From then on it checks silently. If the fingerprint ever \emph{changes}, SSH
refuses to connect and warns loudly --- because that is what an impersonation
attack looks like.

On a home network, the realistic cause of that warning is that you reflashed the
Pi (new key) or your router handed the address to a different device. See
Appendix~\ref{app:trouble0} for the one-line fix --- and note that the fix is to
remove the \emph{specific stale entry}, never to disable checking.
\end{concept}

\section{When \texttt{brewer.local} does not resolve}

\code{.local} names use \textbf{mDNS} --- machines shouting ``I am brewer'' on
the local network, with no central server. It works out of the box on macOS and
most Linux, and on Windows 10+.

When it fails, work down this list.

\subsection{1. Is it a name problem or a network problem?}

Find the Pi's IP from your router's admin page (usually
\code{192.168.1.1} in a browser; look for ``connected devices'' or ``DHCP
leases''). Then:

\begin{shell}
ssh martin@192.168.1.42
\end{shell}

If \textbf{that works}, the Pi is fine and only the name lookup is broken --- fix
it with the \code{\textasciitilde/.ssh/config} entry in
Chapter~\ref{ch:sshkeys}, which lets you type \code{ssh brewer} regardless.

If \textbf{that fails too}, the Pi is not on the network. Continue.

\subsection{2. Scan the network}

\begin{shell}
# macOS / Linux -- find your own address first
ipconfig getifaddr en0        # macOS
hostname -I                   # Linux

# then scan that subnet for open SSH ports
nmap -p 22 --open 192.168.1.0/24
\end{shell}

Install \code{nmap} with \code{brew install nmap} or \code{sudo apt install
nmap}. Any device answering on port 22 is running SSH; if you only have one Pi,
that is it.

\subsection{3. Still nothing}

In order of likelihood:

\begin{center}\small
\begin{tabular}{@{}L{0.42\linewidth}L{0.50\linewidth}@{}}
\toprule
\thead{Cause} & \thead{Check} \\
\midrule
Wi-Fi credentials wrong & Re-flash the card with Imager, retyping the SSID and password carefully. Both are case-sensitive. \\
5\,GHz network selected & A Zero 2 W cannot see 5\,GHz. Use the 2.4\,GHz SSID. \\
Wireless country not set & Imager's \emph{Wireless LAN country} is required; without it the radio stays off. \\
SSH not enabled & Services tab in Imager. Easy to miss. \\
Power & Undervoltage. Try a different supply \emph{and} cable. \\
Card not written & Imager reported an error, or you ejected early. \\
\bottomrule
\end{tabular}
\end{center}

\begin{trap}{Re-flashing is faster than diagnosing}
If the Pi has never appeared, do not spend an hour on it. Re-flash with Imager,
type the Wi-Fi details slowly, and try again. It takes five minutes, and it fixes
the great majority of never-appeared cases.

Diagnosis is worth it once the Pi has worked at least once --- then you know the
hardware is good and something specific changed.
\end{trap}

\section{Look around}

You are logged in. Confirm what you are on:

\begin{shell}
# on the Pi
hostname                              # brewer
whoami                                # martin
cat /proc/device-tree/model; echo     # Raspberry Pi Zero 2 W Rev 1.0
free -h                               # memory: total / used / free
df -h /                               # disk space
uname -m                              # aarch64  <- 64-bit, as intended
ip addr show wlan0 | grep 'inet '     # your IP address
\end{shell}

Two things worth noting now:

\code{free -h} on a Zero 2 W shows about 420\,MB available. That is why
Chapter~\ref{ch:pisetup} adds swap and Chapter~\ref{ch:deploy} builds on your
laptop.

\code{uname -m} must say \code{aarch64}. If it says \code{armv7l} you flashed the
32-bit image --- workable, but re-flash with the 64-bit one now rather than
later.

To leave:

\begin{shell}
exit          # or Ctrl-D
\end{shell}

\section{DHCP, and why the address moves}

Your router hands out addresses on a lease. When the lease expires, or the Pi
reboots, or another device claims it first, the Pi can get a \emph{different}
address.

That is why this book uses \code{brewer.local} rather than a number. If mDNS is
unreliable on your network, the durable fix is a \textbf{static DHCP lease} (also
called an address reservation) in your router: find the Pi in the device list and
tick ``always give this device this address''. Takes a minute, and it is worth
doing before you start filming or demoing anything.

\begin{checkpoint}{You can log in}
\begin{shell}
ssh martin@brewer.local
hostname
exit
\end{shell}

You reached a prompt, it printed \code{brewer}, and you got back out. You typed a
password to do it --- that is what the next chapter fixes.
\end{checkpoint}
